\begin{pastexam}[label={px:nchu-111-is-03}]{中興資工所乙組 111 學年度 第 3 題}
（單選）

(3 pts) What does the following function do for a given Linked List with first node as \textit{head}?
\begin{minted}[fontsize=\small]{c}
void fun1(struct node* head)
{
  if(head == NULL)
    return;

  fun1(head->next);
  printf("%d  ", head->data);
}
\end{minted}
\begin{enumerate}[label=(\alph*).]
\item Prints all nodes of linked lists
\item Prints all nodes of linked list in reverse order
\item Prints alternate nodes of Linked List
\item Prints alternate nodes in reverse order
\end{enumerate}
\end{pastexam}

\begin{pastsoln}
答案是 (b)。

\texttt{fun1} 先遞迴處理「從下一個節點開始的串列」，回來之後才印自己。 以 $1 \to 2 \to 3$ 為例：
\texttt{fun1(1)} 呼叫 \texttt{fun1(2)}，再呼叫 \texttt{fun1(3)}，再呼叫 \texttt{fun1(NULL)} 直接返回；
接著依返回的順序印出：\texttt{fun1(3)} 印 $3$，\texttt{fun1(2)} 印 $2$，\texttt{fun1(1)} 印 $1$。
輸出是 $3\ 2\ 1$。

一般地，對串列長度歸納：空串列什麼都不印；長度 $k \ge 1$ 時先把後面 $k - 1$ 個反向印出（歸納假設），再印
第一個，正好是整條串列的反向。 每個節點都會印到（不會跳過），所以是 (b)，不是 (a) 正向、也不是隔一個印
的 (c)(d)。

遞迴呼叫的深度是串列長度 $n$，用到 $\Theta(n)$ 的堆疊空間（資結書第~6.4~節）；如果先呼叫
\texttt{printf} 再遞迴，印出的就是正向的順序 (a)。
\end{pastsoln}
